\documentclass[10pt,a4paper]{amsart}
\usepackage[margin=19mm]{geometry}
\usepackage{amsmath,amssymb,mathtools}
\usepackage[scr=rsfs,cal=euler]{mathalfa}
\usepackage{xfrac}
\usepackage[unicode,hidelinks,bookmarks=false]{hyperref}
\newcommand{\ie}{i.e.\ }
\newcommand{\cf}{cf.\ }
\newcommand{\resp}{respectively\ }
\newcommand{\Subsection}{\subsection}
\providecommand{\nobreakdash}{\nobreak\hbox{-}}
\setlength{\emergencystretch}{2em}
\multlinegap0pt
\usepackage{bm}
\usepackage{bbm}
\usepackage{dsfont}
\usepackage{xspace}
\usepackage{cancel}
\usepackage{mathtools}
\usepackage{amsthm}
\makeatletter
\@ifundefined{corollary}{\newtheorem{corollary}{Corollary}}{}
\@ifundefined{theorem}{\newtheorem{theorem}{Theorem}}{}
\makeatother
\makeatletter
\expandafter\ifx\csname oneshot\endcsname\relax
\newenvironment{oneshot}[1]{\@begintheorem{#1}{\unskip}}{\@endtheorem}
\fi
\makeatother
\title[Existence and Enumeration of Polynomially Transformed Matric]{Existence and Enumeration of Polynomially Transformed Matrices under Spectral and Nilpotent Constraints}
\author{Shih-Yu Chang}
\date{Source: arXiv:2510.16901v1 (CC BY 4.0)}
\begin{document}
\maketitle

\noindent\emph{Source and licence.} Existence and Enumeration of Polynomially Transformed Matrices under Spectral and Nilpotent Constraints. Version arXiv:2510.16901v1 (CC BY 4.0), distributed under the Creative Commons Attribution 4.0 International licence, as recorded in the arXiv licence metadata for this exact version and shown on the abstract page https://arxiv.org/abs/2510.16901v1. Full rights notice: licenses/arxiv-2510.16901-CC-BY-4.0.txt.\par\smallskip

\noindent\textit{Editorial scope.} Counted unit: the Theorem whose source label is \texttt{thm:Sufficient and checkable conditions for at least K distinct points} in the article (source0.tex lines 1163--1170, source label \texttt{thm:Sufficient and checkable conditions for at least K distinct points}), delivered together with the article's own complete original proof of it (source0.tex lines 1172--1218, one \texttt{proof} environment of 47 lines). No mathematics is written, repaired or re-derived: statement and proof are the article's own text line for line. Required context retained uncounted: cor:Practical counting via squarefree part (lines 1149--1155). Every definition, display and label of those lines is retained unchanged. Omitted from this package and not redistributed: 1--1148, 1156--1162, 1171--1171, 1219--1334. Nothing inside a retained excerpt is omitted or altered: every retained body is delivered exactly as the article's lines are written, so no omission receipt of any kind accompanies this package. Citations: the retained excerpts carry no citation command at all, so no bibliography entry of the article is transcribed and no reference is redistributed. Reference adaptation: none was needed and none was made; every reference inside the delivered unit resolves inside it, so no unresolved reference or citation is produced. Printed slips kept literal: no printed word, symbol, hypothesis or number of the counted statement or of its proof was altered. Layout and class: the article's own class is \texttt{article} with packages including geometry, times, expl3, cite, xcolor, multirow, stackengine, hhline, lipsum, titlesec, wrapfig, enumerate, epsfig, amsmath; the wrapper is the lane's pinned \texttt{amsart} editorial wrapper at 10pt on A4, which restates the theorem-like environments the retained text needs and adds the article's own macro definitions that the retained text needs (none), with extra wrapper packages (bm, bbm, dsfont, xspace, cancel, mathtools). The delivered numbering is the unit's own. Assets: no retained span contains a figure environment, a graphics-inclusion command, a PDF-page inclusion, a table or a TikZ picture; no image file is copied into this package, the package declares no asset dependency and no image file is redistributed with it. Licence: version arXiv:2510.16901v1 of this article is distributed under the Creative Commons Attribution 4.0 International licence, as recorded in the arXiv licence metadata for this exact version and shown on the abstract page https://arxiv.org/abs/2510.16901v1. A full rights notice is delivered at licenses/arxiv-2510.16901-CC-BY-4.0.txt. Characters: every retained excerpt is pure ASCII, so the delivered engine, XeLaTeX with its default Latin Modern text font, reports no missing character.
\begin{corollary}[Practical counting via squarefree part]\label{cor:Practical counting via squarefree part}
Let \(h_{\mathrm{sf}}(x)\) denote the squarefree part of \(h(x)\),
\[
h_{\mathrm{sf}}(x)=\frac{h(x)}{\gcd(h(x),h'(x))}.
\]
Then the number of distinct points \(x\) with \(f(x)\neq0\) and \(f',\dots,f^{(\acute{m}-1)}\) vanishing equals \(\deg h_{\mathrm{sf}}\).
\end{corollary}

\begin{theorem}[Sufficient and checkable conditions for at least \(K\) distinct points]\label{thm:Sufficient and checkable conditions for at least K distinct points}
With notation as above, the following are sufficient (and easily checkable) conditions that guarantee \(\#\mathcal{Z}\ge K\), where $\mathcal{Z}$ is the number of distinct roots for the polynomial $f$:
\begin{enumerate}
  \item \(\deg h_{\mathrm{sf}}\ge K\).
  \item Equivalently, \(\deg h \ge K\) and \(\gcd(h,h')\) has degree at most \(\deg h-K\) (so that after removing repeated factors at least \(K\) distinct linear factors remain).
  \item Computationally: compute \(g=\gcd(f',\dots,f^{(\acute{m}-1)})\), \(d=\gcd(f,g)\), form \(h=g/d\), then compute the squarefree part \(h_{\mathrm{sf}}=h/\gcd(h,h')\). If \(\deg h_{\mathrm{sf}}\ge K\) then there are at least \(K\) distinct points in \(\mathcal{Z}\).
\end{enumerate}
\end{theorem}

\begin{proof}

We prove the three items in order, reducing each to Corollary~\ref{cor:Practical counting via squarefree part}.

\medskip\noindent\textbf{(1)} By construction (and the discussion preceding the theorem) the polynomial \(h\) has exactly the same roots (counted with multiplicity) as the set of points \(z\) for which
\[
f'(x)=f''(x)=\cdots=f^{(\acute{m}-1)}(x)=0
\]
but without the roots where additionally \(f(x)=0\) (those have been removed when forming \(h\); see item (3) for the explicit construction). Therefore the number of \emph{distinct} points \(z\) with \(f(x)\neq0\) and \(f',\dots,f^{(\acute{m}-1)}\) vanishing is exactly the number of distinct roots of \(h\). By Corollary~\ref{cor:Practical counting via squarefree part} this number equals \(\deg h_{\mathrm{sf}}\). Hence if \(\deg h_{\mathrm{sf}}\ge K\) then \(\#\mathcal Z\ge K\), proving item 1.

\medskip\noindent\textbf{(2)} The equality
\[
h_{\mathrm{sf}}=\frac{h}{\gcd(h,h')}
\]
implies
\[
\deg h_{\mathrm{sf}}=\deg h-\deg\big(\gcd(h,h')\big).
\]
Thus \(\deg h_{\mathrm{sf}}\ge K\) is equivalent to
\[
\deg h-\deg\big(\gcd(h,h')\big)\ge K,
\]
i.e.
\[
\deg\big(\gcd(h,h')\big)\le \deg h-K.
\]
Interpreting \(\gcd(h,h')\) as the polynomial collecting all repeated prime factors of \(h\), this condition means that after removing repeated factors from \(h\) at least \(K\) distinct linear factors remain. This proves the equivalence asserted in item 2.

\medskip\noindent\textbf{(3)} We now justify the computational recipe. Let
\[
g:=\gcd\big(f',f'',\dots,f^{(\acute{m}-1)}\big).
\]
By definition every root of \(g\) is a point where all derivatives \(f',\dots,f^{(\acute{m}-1)}\) vanish; conversely every point where those derivatives vanish is a root of \(g\). Some of those roots might also satisfy \(f(x)=0\); to exclude those we form
\[
d:=\gcd(f,g),
\]
so that \(d\) captures exactly the common factors corresponding to points where \(f=0\) and the derivatives vanish simultaneously. Then
\[
h:=\frac{g}{d}
\]
is a polynomial whose roots are precisely the points where the derivatives vanish but \(f\neq0\), counted with their multiplicities inherited from \(g\). Finally compute the squarefree part
\[
h_{\mathrm{sf}}=\frac{h}{\gcd(h,h')},
\]
which has the same distinct roots as \(h\) but with multiplicity \(1\). By Corollary~\ref{cor:Practical counting via squarefree part} the number of distinct such points equals \(\deg h_{\mathrm{sf}}\). Hence \(\deg h_{\mathrm{sf}}\ge K\) implies \(\#\mathcal Z\ge K\), which is exactly the computational criterion stated in item 3. This completes the proof of Theorem~\ref{thm:Sufficient and checkable conditions for at least K distinct points}.
%
\end{proof}

\end{document}
